% =============================================================================
\chapter{Introduction}
% =============================================================================

\section{Objet du document}

Ce \textbf{cahier d'analyse et de conception (CAC)} constitue la \emph{documentation de référence}
du backend \yow. Il décrit l'ensemble du système : contexte métier, acteurs, cas d'utilisation,
architecture logicielle, modèles de données, intégrations externes, sécurité et déploiement.

Ce document s'adresse aux développeurs, architectes, testeurs et parties prenantes techniques
souhaitant comprendre, maintenir ou faire évoluer la plateforme.

\section{Présentation du projet}

\yow\ est une plateforme \textbf{SaaS multi-tenant} dédiée aux \textbf{artistes visuels} et aux
\textbf{collectionneurs}. Elle permet :

\begin{itemize}[noitemsep]
  \item aux artistes de créer un espace personnel (tenant isolé), publier des œuvres, gérer une boutique en ligne, organiser des événements et recevoir des paiements ;
  \item aux acheteurs de découvrir des artistes, acheter des œuvres, réserver des billets et interagir (likes, commentaires, chat) ;
  \item aux administrateurs de valider les inscriptions artistes, superviser les tenants et consulter les statistiques globales.
\end{itemize}

\section{Périmètre}

\begin{table}[h]
  \centering
  \caption{Périmètre fonctionnel du backend}
  \label{tab:perimetre}
  \begin{tabular}{@{}lp{10cm}@{}}
    \toprule
    \textbf{Inclus} & Authentification JWT, multi-tenancy, œuvres, boutique, événements, paiements CamPay, wallet, abonnements, chat WebSocket, notifications, recherche, administration \\
    \midrule
    \textbf{Exclus} & Frontend Next.js, interface utilisateur Kernel RT-Comops, infrastructure réseau (load balancers, CDN) \\
    \bottomrule
  \end{tabular}
\end{table}

\section{Stack technique}

\begin{table}[h]
  \centering
  \caption{Technologies utilisées}
  \label{tab:stack}
  \begin{tabular}{@{}ll@{}}
    \toprule
    \textbf{Composant} & \textbf{Technologie / Version} \\
    \midrule
    Langage            & Java 17 \\
    Framework          & Spring Boot 4.0.0 (Spring Framework 7.0) \\
    Sécurité           & Spring Security 7.0, OAuth2 Resource Server, JWT RS256 \\
    Persistance        & Spring Data JPA, Hibernate 6, PostgreSQL 18 \\
    Multi-tenancy      & Schema-per-tenant (Hibernate) \\
    Migrations         & Flyway 11.0.0 \\
    Paiements          & CamPay (Mobile Money MTN, Orange, MoMo Pay) \\
    Messagerie         & WebSocket STOMP (Spring) \\
    Événements         & Apache Kafka (optionnel, sync Kernel) \\
    Documentation API  & OpenAPI 3 / SpringDoc 3.0.0 \\
    E-mail             & Spring Boot Mail \\
    Build              & Maven 3.9+ \\
    \bottomrule
  \end{tabular}
\end{table}

\section{Conventions documentaires}

\begin{itemize}[noitemsep]
  \item Les diagrammes UML sont réalisés en \TikZ\ pour une compilation LaTeX autonome.
  \item Les chemins de packages suivent la convention \texttt{com.yowpainter.*}.
  \item Les endpoints API sont préfixés par \texttt{/api}.
  \item Les identifiants UUID sont notés \texttt{UUID} sauf mention contraire.
\end{itemize}

\section{Références}

\begin{itemize}[noitemsep]
  \item \texttt{README.md} — Vue d'ensemble et installation
  \item \texttt{docs/KERNEL\_INTEGRATION.md} — Intégration Kernel RT-Comops
  \item \texttt{docs/ARTIST\_APPROVAL\_WORKFLOW.md} — Workflow validation artiste
  \item \texttt{docs/CONFIGURATION\_KERNEL\_YOWPAINTER.md} — Configuration locale
  \item Swagger UI : \url{http://localhost:8090/swagger-ui.html}
\end{itemize}
